This paper presents the evolution of the \textit{FireSafe} system towards a distributed Internet of Things (IoT) architecture, designed for critical environmental monitoring and early hazard detection. Unlike its predecessor, this version implements a network topology where multiple ESP32-based peripheral nodes perform edge computing, managing task concurrency through the FreeRTOS real-time operating system \citep{valvano2014volume1}. 

The infrastructure is supported by the MQTT asynchronous messaging protocol under a publish/subscribe model, ensuring efficient and scalable communication between nodes and the central server on a Raspberry Pi. The system integrates a Django-developed backend for data persistence and a React-based interactive dashboard for real-time visualization and automatic alert management. The results demonstrate that the circuit analysis proposed by \citet{boylestad2009} and the implementation of real-time interfacing \citep{valvano2014volume2} allow for an immediate response to thermal fluctuations or the presence of combustible gases, validating the feasibility of low-cost industrial solutions.

\textbf{Keywords:} Distributed IoT, MQTT, FreeRTOS, environmental monitoring, real-time systems.
